\section{AHSME 1977}
        \begin{enumerate}
        	\item (\textit{AHSME 1977}) If $y = 2x$ and $z = 2y$, then $x + y + z$ equals 


 $\text{(A)}\ x \qquad  \text{(B)}\ 3x \qquad  \text{(C)}\ 5x \qquad  \text{(D)}\ 7x \qquad  \text{(E)}\ 9x$


 

	\item (\textit{AHSME 1977}) Which one of the following statements is false? All equilateral triangles are


 $\textbf{(A)}\ \text{ equiangular}\qquad \textbf{(B)}\ \text{isosceles}\qquad \textbf{(C)}\ \text{regular polygons }\qquad\\ \textbf{(D)}\ \text{congruent to each other}\qquad \textbf{(E)}\ \text{similar to each other}$


 

	\item (\textit{AHSME 1977}) A man has $2.73 in pennies, nickels, dimes, quarters and half dollars. If he has an equal number of coins of each kind, then the total number of coins he has is 


 $\text{(A)}\ 3 \qquad  \text{(B)}\ 5 \qquad  \text{(C)}\ 9 \qquad  \text{(D)}\ 10 \qquad  \text{(E)}\ 15$


 

	\item (\textit{AHSME 1977}) 
                    \begin{figure}[H]
                        \centering
                        \includegraphics[width=0.5\textwidth]{Problem/AHSME_I_1977_4_Problem_Figure_1.png}
                    \end{figure}
                    


 In triangle $ABC, AB=AC$ and $\measuredangle A=80^\circ$. If points $D, E$, and $F$ lie on sides $BC, AC$ and $AB$, respectively, and $CE=CD$ and $BF=BD$, then $\measuredangle EDF$ equals


 $\textbf{(A) }30^\circ\qquad \textbf{(B) }40^\circ\qquad \textbf{(C) }50^\circ\qquad \textbf{(D) }65^\circ\qquad \textbf{(E) }\text{none of these}$


 

	\item (\textit{AHSME 1977}) The set of all points $P$ such that the sum of the (undirected) distances from $P$ to two fixed points $A$ and $B$ equals the distance between $A$ and $B$ is


 $\textbf{(A) }\text{the line segment from }A\text{ to }B\qquad \textbf{(B) }\text{the line passing through }A\text{ and }B\qquad\\ \textbf{(C) }\text{the perpendicular bisector of the line segment from }A\text{ to }B\qquad\\ \textbf{(D) }\text{an ellipse having positive area}\qquad \textbf{(E) }\text{a parabola}$


 

	\item (\textit{AHSME 1977}) If $x, y$ and $2x + \frac{y}{2}$ are not zero, then
$\left( 2x + \frac{y}{2} \right)^{-1} \left[(2x)^{-1} + \left( \frac{y}{2} \right)^{-1} \right]$
equals


 $\textbf{(A) }1\qquad \textbf{(B) }xy^{-1}\qquad \textbf{(C) }x^{-1}y\qquad \textbf{(D) }(xy)^{-1}\qquad  \textbf{(E) }\text{none of these}$


 

	\item (\textit{AHSME 1977}) If $t = \frac{1}{1 - \sqrt[4]{2}}$, then $t$ equals 


 $\text{(A)}\ (1-\sqrt[4]{2})(2-\sqrt{2})\qquad \text{(B)}\ (1-\sqrt[4]{2})(1+\sqrt{2})\qquad \text{(C)}\ (1+\sqrt[4]{2})(1-\sqrt{2}) \qquad \\ \text{(D)}\ (1+\sqrt[4]{2})(1+\sqrt{2})\qquad \text{(E)}-(1+\sqrt[4]{2})(1+\sqrt{2})$


 



 

	\item (\textit{AHSME 1977}) For every triple $(a,b,c)$ of non-zero real numbers, form the number $\frac{a}{|a|}+\frac{b}{|b|}+\frac{c}{|c|}+\frac{abc}{|abc|}$. 
The set of all numbers formed is


 $\textbf{(A)}\ {0} \qquad \textbf{(B)}\ \{-4,0,4\} \qquad \textbf{(C)}\ \{-4,-2,0,2,4\} \qquad \textbf{(D)}\ \{-4,-2,2,4\}\qquad \textbf{(E)}\ \text{none of these}$


 

	\item (\textit{AHSME 1977}) 
                    \begin{figure}[H]
                        \centering
                        \includegraphics[width=0.5\textwidth]{Problem/AHSME_I_1977_9_Problem_Figure_1.png}
                    \end{figure}
                    


 

In the adjoining figure $\measuredangle E=40^\circ$ and arc $AB$, arc $BC$, and arc $CD$ all have equal length. Find the measure of $\measuredangle ACD$.


 $\textbf{(A) }10^\circ\qquad \textbf{(B) }15^\circ\qquad \textbf{(C) }20^\circ\qquad \textbf{(D) }\left(\frac{45}{2}\right)^\circ\qquad \textbf{(E) }30^\circ$


 

	\item (\textit{AHSME 1977}) If $(3x-1)^7 = a_7x^7 + a_6x^6 + \cdots + a_0$, then $a_7 + a_6 + \cdots + a_0$ equals 


 $\text{(A)}\ 0 \qquad  \text{(B)}\ 1 \qquad  \text{(C)}\ 64 \qquad  \text{(D)}\ -64 \qquad  \text{(E)}\ 128$


 

	\item (\textit{AHSME 1977}) For each real number $x$, let $\textbf{[}x\textbf{]}$ be the largest integer not exceeding $x$ 
(i.e., the integer $n$ such that $n\le x<n+1$). Which of the following statements is (are) true?


 $\textbf{I. [}x+1\textbf{]}=\textbf{[}x\textbf{]}+1\text{ for all }x \\ \textbf{II. [}x+y\textbf{]}=\textbf{[}x\textbf{]}+\textbf{[}y\textbf{]}\text{ for all }x\text{ and }y \\ \textbf{III. [}xy\textbf{]}=\textbf{[}x\textbf{]}\textbf{[}y\textbf{]}\text{ for all }x\text{ and }y$


 

$\textbf{(A) }\text{none}\qquad \textbf{(B) }\textbf{I }\text{only}\qquad \textbf{(C) }\textbf{I}\text{ and }\textbf{II}\text{ only}\qquad \textbf{(D) }\textbf{III }\text{only}\qquad \textbf{(E) }\text{all}$


 

	\item (\textit{AHSME 1977}) Al's age is $16$ more than the sum of Bob's age and Carl's age, and the square of Al's age is $1632$ more than the square of the sum of 
Bob's age and Carl's age. What is the sum of the ages of Al, Bob, and Carl?     


 $\text{(A)}\ 64 \qquad  \text{(B)}\ 94 \qquad  \text{(C)}\ 96 \qquad  \text{(D)}\ 102 \qquad  \text{(E)}\ 140$


 

	\item (\textit{AHSME 1977}) If $a_1,a_2,a_3,\dots$ is a sequence of positive numbers such that $a_{n+2}=a_na_{n+1}$ for all positive integers $n$, 
then the sequence $a_1,a_2,a_3,\dots$ is a geometric progression


 $\textbf{(A) }\text{for all positive values of }a_1\text{ and }a_2\qquad\\ \textbf{(B) }\text{if and only if }a_1=a_2\qquad\\ \textbf{(C) }\text{if and only if }a_1=1\qquad\\ \textbf{(D) }\text{if and only if }a_2=1\qquad\\ \textbf{(E) }\text{if and only if }a_1=a_2=1$


 

	\item (\textit{AHSME 1977}) How many pairs $(m,n)$ of integers satisfy the equation $m+n=mn$?


 $\textbf{(A) }1\qquad \textbf{(B) }2\qquad \textbf{(C) }3\qquad \textbf{(D) }4\qquad  \textbf{(E) }\text{more than }4$


 

	\item (\textit{AHSME 1977}) 
                    \begin{figure}[H]
                        \centering
                        \includegraphics[width=0.5\textwidth]{Problem/AHSME_I_1977_15_Problem_Figure_1.png}
                    \end{figure}
                    


 Each of the three circles in the adjoining figure is externally tangent to the other two, and each side of the triangle is tangent to two of the circles. If each circle has radius three, then the perimeter of the triangle is


 $\textbf{(A) }36+9\sqrt{2}\qquad \textbf{(B) }36+6\sqrt{3}\qquad \textbf{(C) }36+9\sqrt{3}\qquad \textbf{(D) }18+18\sqrt{3}\qquad \textbf{(E) }45$


 

	\item (\textit{AHSME 1977}) If $i^2 = -1$, then the sum 
 \[\cos{45^\circ} + i\cos{135^\circ} + \cdots + i^n\cos{(45 + 90n)^\circ}  + \cdots + i^{40}\cos{3645^\circ}\]
equals 


 $\text{(A)}\ \frac{\sqrt{2}}{2} \qquad  \text{(B)}\ -10i\sqrt{2} \qquad  \text{(C)}\ \frac{21\sqrt{2}}{2} \qquad\\ \text{(D)}\ \frac{\sqrt{2}}{2}(21 - 20i) \qquad  \text{(E)}\ \frac{\sqrt{2}}{2}(21 + 20i)$


 

	\item (\textit{AHSME 1977}) Three fair dice are tossed at random (i.e., all faces have the same probability of coming up). 
What is the probability that the three numbers turned up can be arranged to form an arithmetic progression with common difference one?


 $\textbf{(A) }\frac{1}{6}\qquad \textbf{(B) }\frac{1}{9}\qquad \textbf{(C) }\frac{1}{27}\qquad \textbf{(D) }\frac{1}{54}\qquad \textbf{(E) }\frac{7}{36}$


 

	\item (\textit{AHSME 1977}) If $y=(\log_23)(\log_34)\cdots(\log_n[n+1])\cdots(\log_{31}32)$ then 


 $\textbf{(A) }4<y<5\qquad \textbf{(B) }y=5\qquad \textbf{(C) }5<y<6\qquad \textbf{(D) }y=6\qquad \\ \textbf{(E) }6<y<7$


 

	\item (\textit{AHSME 1977}) Let $E$ be the point of intersection of the diagonals of convex quadrilateral $ABCD$, and let $P,Q,R$, and $S$ be the centers of the circles 
circumscribing triangles $ABE, BCE, CDE$, and $ADE$, respectively. Then


 $\textbf{(A) }PQRS\text{ is a parallelogram}\\ \textbf{(B) }PQRS\text{ is a parallelogram if an only if }ABCD\text{ is a rhombus}\\ \textbf{(C) }PQRS\text{ is a parallelogram if an only if }ABCD\text{ is a rectangle}\\ \textbf{(D) }PQRS\text{ is a parallelogram if an only if }ABCD\text{ is a parallelogram}\\ \textbf{(E) }\text{none of the above are true}$


 

	\item (\textit{AHSME 1977}) $\begin{tabular}{ccccccccccccc}& & & & & & C & & & & & &\\  & & & & & C & O & C & & & & &\\  & & & & C & O & N & O & C & & & &\\  & & & C & O & N & T & N & O & C & & &\\  & & C & O & N & T & E & T & N & O & C & &\\  & C & O & N & T & E & S & E & T & N & O & C &\\  C & O & N & T & E & S & T & S & E & T & N & O & C \end{tabular}$


 

For how many paths consisting of a sequence of horizontal and/or vertical line segments, with each segment connecting a pair of adjacent letters in the diagram above, is the word CONTEST spelled out as the path is traversed from beginning to end?


 $\textbf{(A) }63\qquad \textbf{(B) }128\qquad \textbf{(C) }129\qquad \textbf{(D) }255\qquad  \textbf{(E) }\text{none of these}$


 

	\item (\textit{AHSME 1977}) For how many values of the coefficient a do the equations 
 \begin{align*}x^2+ax+1=0 \\ x^2-x-a=0\end{align*}
 have a common real solution?


 $\textbf{(A)}\ 0 \qquad \textbf{(B)}\ 1 \qquad \textbf{(C)}\ 2 \qquad \textbf{(D)}\ 3 \qquad \textbf{(E)}\ \infty$


 

	\item (\textit{AHSME 1977}) If $f(x)$ is a real valued function of the real variable $x$, and $f(x)$ is not identically zero, 
and for all $a$ and $b$ $f(a+b)+f(a-b)=2f(a)+2f(b)$, then for all $x$ and $y$


 $\textbf{(A) }f(0)=1\qquad \textbf{(B) }f(-x)=-f(x)\qquad \textbf{(C) }f(-x)=f(x)\qquad \\ \textbf{(D) }f(x+y)=f(x)+f(y) \qquad \\ \textbf{(E) }\text{there is a positive real number }T\text{ such that }f(x+T)=f(x)$


 

	\item (\textit{AHSME 1977}) If the solutions of the equation $x^2+px+q=0$ are the cubes of the solutions of the equation $x^2+mx+n=0$, then


 $\textbf{(A) }p=m^3+3mn\qquad \textbf{(B) }p=m^3-3mn\qquad \textbf{(C) }p+q=m^3\qquad\\ \textbf{(D) }\left(\frac{m}{n}\right)^2=\frac{p}{q}\qquad  \textbf{(E) }\text{none of these}$


 

	\item (\textit{AHSME 1977}) Find the sum $\frac{1}{1(3)}+\frac{1}{3(5)}+\dots+\frac{1}{(2n-1)(2n+1)}+\dots+\frac{1}{255(257)}$.


 $\textbf{(A) }\frac{127}{255}\qquad \textbf{(B) }\frac{128}{255}\qquad \textbf{(C) }\frac{1}{2}\qquad \textbf{(D) }\frac{128}{257}\qquad \textbf{(E) }\frac{129}{257}$


 

	\item (\textit{AHSME 1977}) Determine the largest positive integer n such that $1005!$ is divisible by $10^n$.


 $\textbf{(A) }102\qquad \textbf{(B) }112\qquad \textbf{(C) }249\qquad \textbf{(D) }502\qquad  \textbf{(E) }\text{none of these}$


 

	\item (\textit{AHSME 1977}) Let $a,b,c$, and $d$ be the lengths of sides $MN,NP,PQ$, and $QM$, respectively, of quadrilateral $MNPQ$. If $A$ is the area of $MNPQ$, then


 $\textbf{(A) }A=\left(\frac{a+c}{2}\right)\left(\frac{b+d}{2}\right)\text{ if and only if }MNPQ\text{ is convex}\\ \textbf{(B) }A=\left(\frac{a+c}{2}\right)\left(\frac{b+d}{2}\right)\text{ if and only if }MNPQ\text{ is a rectangle}\\ \textbf{(C) }A\le\left(\frac{a+c}{2}\right)\left(\frac{b+d}{2}\right)\text{ if and only if }MNPQ\text{ is a rectangle}\\ \textbf{(D) }A\le\left(\frac{a+c}{2}\right)\left(\frac{b+d}{2}\right)\text{ if and only if }MNPQ\text{ is a parallelogram}\\ \textbf{(E) }A\ge\left(\frac{a+c}{2}\right)\left(\frac{b+d}{2}\right)\text{ if and only if }MNPQ\text{ is a parallelogram}$


 

	\item (\textit{AHSME 1977}) There are two spherical balls of different sizes lying in two corners of a rectangular room, each touching two walls and the floor. 
If there is a point on each ball which is $5$ inches from each wall which that ball touches and $10$ inches from the floor, 
then the sum of the diameters of the balls is


 $\textbf{(A) }20\text{ inches}\qquad \textbf{(B) }30\text{ inches}\qquad \textbf{(C) }40\text{ inches}\qquad \textbf{(D) }60\text{ inches}\qquad \\ \textbf{(E) }\text{not determined by the given information}$


 

	\item (\textit{AHSME 1977}) Let $g(x)=x^5+x^4+x^3+x^2+x+1$. What is the remainder when the polynomial $g(x^{12})$ is divided by the polynomial $g(x)$?


 $\textbf{(A) }6\qquad \textbf{(B) }5-x\qquad \textbf{(C) }4-x+x^2\qquad \textbf{(D) }3-x+x^2-x^3\qquad \\ \textbf{(E) }2-x+x^2-x^3+x^4$


 

	\item (\textit{AHSME 1977}) Find the smallest integer $n$ such that $(x^2+y^2+z^2)^2\le n(x^4+y^4+z^4)$ for all real numbers $x,y$, and $z$.


 $\textbf{(A) }2\qquad \textbf{(B) }3\qquad \textbf{(C) }4\qquad \textbf{(D) }6\qquad  \textbf{(E) }\text{There is no such integer n}$


 

	\item (\textit{AHSME 1977}) 
                    \begin{figure}[H]
                        \centering
                        \includegraphics[width=0.5\textwidth]{Problem/AHSME_I_1977_30_Problem_Figure_1.png}
                    \end{figure}
                    


 If $a,b$, and $d$ are the lengths of a side, a shortest diagonal and a longest diagonal, respectively, 
of a regular nonagon (see adjoining figure), then


 $\textbf{(A) }d=a+b\qquad \textbf{(B) }d^2=a^2+b^2\qquad \textbf{(C) }d^2=a^2+ab+b^2\qquad\\ \textbf{(D) }b=\frac{a+d}{2}\qquad  \textbf{(E) }b^2=ad$


 

\end{enumerate}